\section{Two phases and one final clipping operation}
\label{sec:revisit-completion}
The two phases are a geometric regularization ramp and an accuracy phase
at its final regularizer. They share both numerical vectors and all graph
state. Neither phase makes a restricted-system solve or invokes SDD.

\begin{lemma}[An explicit accuracy-phase energy bound]
\label{lem:revisit-freeze-energy}
At any checkpoint of the cooling recurrence, with
$r=\lambda/\alpha$ and $\vu=\vx_r^*$, the Euclidean accelerated energy
satisfies
\begin{equation}
 E=J_\lambda(\vx)-J_\lambda(\vu)
                           +\tfrac\mu2\|\vz-\vu\|_2^2\leq9\lambda.
 \label{eq:revisit-freeze-energy}
\end{equation}
\end{lemma}
\begin{proof}
The response bound and $\mQ\succeq\alpha\mI$ give
\[
 f(\vx)-f(\vt)=\tfrac12\|\vx-\vt\|_{\mQ}^2\leq2\lambda,
 \qquad \tfrac\mu2\|\vz-\vt\|_2^2\leq2\lambda.
\]
Also $\vomega^\trans\vx\leq2$, while
$J_\lambda(\vu)\geq f(\vt)$, so the objective part of $E$ is at most
$4\lambda$. The standard RPPR bias gives
$\vzero\leq\vt-\vu\leq r\vomega$ and
$\vomega^\trans(\vt-\vu)\leq1$. Thus
$\|\vt-\vu\|_2^2\leq r$. The squared triangle inequality yields
\[
 \tfrac\mu2\|\vz-\vu\|_2^2
 \leq\mu\|\vz-\vt\|_2^2+\mu\|\vt-\vu\|_2^2
 \leq4\lambda+\mu r\leq5\lambda.
\]
\end{proof}

\begin{algorithm}[t]
\caption{Two-phase orthant acceleration for semantic PPR}
\label{alg:revisit-ppr}
\begin{algorithmic}[1]
\REQUIRE Seed $v$, local graph access, $\alpha\in(0,1]$, $0<\epsppr<1$.
\STATE Read $d_v$; put $\rho=\epsppr/2$.
\IF{$\rho d_v\geq1$}
 \RETURN zero.
\ENDIF
\IF{$\alpha=1$}
 \RETURN the single coordinate $\widehat{\vpi}=(1-\rho d_v)\eunit v$.
\ENDIF
\STATE Choose \cref{eq:revisit-parameters}; set
       $r=1/d_v$ and $\vx=\vz=\vzero$.
\WHILE{$r>\rho$}
 \STATE Apply \cref{eq:revisit-recurrence} with $\lambda=\alpha r$.
 \STATE $r\gets\max\{\rho,\eta r\}$.
\ENDWHILE
\STATE Set $a=9\alpha\rho$ and $\tau=\alpha\epsppr^2/32$.
\WHILE{$a>\tau$}
 \STATE Apply the same recurrence with $\lambda=\alpha\rho$.
 \STATE $a\gets\chi a$.
\ENDWHILE
\STATE Materialize the historical primal support once and form
       $\widehat\vx=[\vx-(\epsppr/4)\vomega]_+$.
\RETURN $\widehat{\vpi}=\mD^{1/2}\widehat\vx$ as a sparse list.
\end{algorithmic}
\end{algorithm}

\begin{theorem}[Direct deterministic semantic PPR]
\label{thm:revisit-semantic}
In the stated exact-real local model, \cref{alg:revisit-ppr} returns
\begin{equation}
 \vzero\leq\widehat\vx\leq\vx_\rho^*\leq\vx_0^*,\qquad
 \vol(\supp\widehat\vx)\leq2/\epsppr,
 \qquad
 \|\mD^{-1/2}(\widehat\vx-\vx_0^*)\|_\infty\leq\epsppr.
 \label{eq:revisit-semantic}
\end{equation}
If $N$ is the number of accuracy-phase steps and $K$ the total number of
steps, then
\begin{equation}
 N\leq1+\theta^{-1}\log(144/\epsppr),\qquad
 \sum_{k<K}\vol(\supp\vz_{k+1})
       \leq\frac{44}{\rho}(N+2/\theta).
 \label{eq:revisit-explicit-work}
\end{equation}
There is a deterministic sparse implementation with total work
\begin{equation}
 O\!\left(
 \frac{\log(144/\epsppr)}{\epsppr\sqrt\alpha}
 \log\!\left(2+
       \frac{\log(144/\epsppr)}{\epsppr\sqrt\alpha}\right)\right).
 \label{eq:revisit-total-work}
\end{equation}
Peak state is
$O(1+\log(144/\epsppr)/(\epsppr\sqrt\alpha))$ words.
The one-sided vector statement does not by itself assert a nonnegative
PageRank residual or an ACL certificate.
\end{theorem}
\begin{proof}[Accuracy and kinetic work]
At the start of the frozen phase, \cref{lem:revisit-freeze-energy} supplies
the computable upper bound $a=9\alpha\rho$. With $\lambda$ fixed,
Euclidean orthant projection and
\cref{eq:revisit-imported-comparison} give $E^+\leq\chi E$.
Thus the scalar stopping rule proves an objective gap at most $\tau$.
Strong convexity gives
\[
 \|\vx-\vx_\rho^*\|_2\leq\sqrt{2\tau/\alpha}=\epsppr/4.
\]
Since $d_i\geq1$, this is also a degree-normalized coordinate error at
most $\epsppr/4$. The downward clipping gives
\[
 \vzero\leq\widehat\vx\leq\vx_\rho^*,\qquad
 \vzero\leq\vx_\rho^*-\widehat\vx\leq(\epsppr/2)\vomega.
\]
Add the RPPR bias at $\rho=\epsppr/2$ and use support containment.
The zero-return branch is correct because then $\vx_\rho^*=\vzero$;
the $\alpha=1$ branch is the exact RPPR vector and has semantic error $\rho$.

The ratio $a/\tau$ at the freeze is $144/\epsppr$ and
$\chi^j\leq\exp(-\theta j)$, giving the bound on $N$.
Before the first clipped-to-$\rho$ parameter, the inverse regularizers
form a geometric series whose sum is less than
$1/((1-\eta)\rho)=2/(\theta\rho)$.
The frozen phase contributes $N/\rho$. Apply
\cref{thm:revisit-volume}. The local implementation and its other charges
are supplied below.
\end{proof}

\subsection{The same handoff also solves the RPPR objective problem}
\begin{corollary}[Direct RPPR objective completion]
\label{cor:revisit-rppr}
For an independently specified $\rho>0$ and $0<\epsobj<1$, use the
same ramp to $\rho$. Choose a dyadic $\delta$ by halving $1/2$ until
$\delta^2\leq\epsobj\rho/2$, and freeze until the scalar energy bound
is at most $\alpha\delta^2/2$. Return
$\widehat\vx=[\vx-\delta\vomega]_+$. In the nonzero regime
$\rho<1/d_v$, this returns
\[
 0\leq\widehat\vx\leq\vx_\rho^*,\qquad
 \vol(\supp\widehat\vx)\leq1/\rho,\qquad
 F_\rho(\widehat\vx)-F_\rho(\vx_\rho^*)\leq\epsobj.
\]
The frozen phase has at most
$1+\theta^{-1}\log(144/\epsobj)$ steps, and fully charged exact-real
work is $\widetilde O(1/(\rho\sqrt\alpha))$ with logarithmic dependence
on $1/\epsobj$. The same constant-output branches apply, using the
requested $\rho$.
\end{corollary}
\begin{proof}
The chosen grid obeys
$\epsobj\rho/8<\delta^2\leq\epsobj\rho/2$.
Strong convexity and the frozen energy bound give
$\|\vx-\vx_\rho^*\|_2\leq\delta$.
Clipping therefore yields
$0\leq\vx_\rho^*-\widehat\vx\leq2\delta\vomega$ and zero error
outside $\supp\vx_\rho^*$.
The KKT gradient of $J_{\alpha\rho}$ vanishes on that support, so
\[
 F_\rho(\widehat\vx)-F_\rho(\vx_\rho^*)
 =\tfrac12\|\widehat\vx-\vx_\rho^*\|_{\mQ}^2
 \leq\frac{2\delta^2}{\rho}\leq\epsobj.
\]
The actual handoff energy is $9\alpha\rho$, whose ratio to the target
is $18\rho/\delta^2<144/\epsobj$. The same kinetic and reporter charges
apply to this alternative stopping rule and terminal clipping.
\end{proof}

\subsection{Optional positive residual completion}
\begin{corollary}[The older AESP--CD KKT diagnostic]
\label{cor:revisit-kkt}
For a nonnegative point, let $\bm{\kappa}_\rho$ be the coordinatewise
minimum-magnitude element of $\partial F_\rho$ and define
$R_{\mathrm{kkt},\rho}(\vx)
=\|\mD^{-1/2}\bm{\kappa}_\rho(\vx)\|_\infty$.
If the frozen gap is at most $\alpha\delta^2/2$, the clipped output
satisfies $R_{\mathrm{kkt},\rho}(\widehat\vx)\leq2\delta$.
Thus taking $\delta=\alpha\epskkt/2$ gives the note-scoped target
$R_{\mathrm{kkt},\rho}(\widehat\vx)/\alpha\leq\epskkt$
with only logarithmic dependence on $1/\epskkt$ in the iteration count.
This concerns the new recurrence, not the unchanged AESP--CD method.
\end{corollary}
\begin{proof}
As above, $\ve=\vx_\rho^*-\widehat\vx$ lies in
$[0,2\delta\vomega]$, so $|\mQ\ve|\leq2\delta\vomega$.
At a positive output coordinate the optimal RPPR gradient is zero, and
the KKT vector equals $-(\mQ\ve)_i$.
At a zero output coordinate the PPR residual is nonnegative; its KKT
violation is the positive excess of that residual over
$\alpha\rho\omega_i$. The optimum's residual is at most that threshold,
so this violation is also bounded by $2\delta\omega_i$.
\end{proof}

\begin{corollary}[A positive residual certificate by tighter final clipping]
\label{cor:revisit-positive-residual}
Let the frozen objective gap be at most $\alpha\delta^2/2$ and suppose
$0<\delta\leq\alpha\rho/2$. The same final clipping
$\widehat\vx=[\vx-\delta\vomega]_+$ then gives
\begin{equation}
 0\leq\vb-\mQ\widehat\vx\leq2\alpha\rho\vomega.
 \label{eq:revisit-positive-residual}
\end{equation}
For $\rho=\epsppr/2$, this is a nonnegative PPR residual certificate
at semantic accuracy $\epsppr$. It can be obtained with the same
$\widetilde O(1/(\epsppr\sqrt\alpha))$ work by freezing to
$\alpha^3\epsppr^2/32$ and clipping by $\alpha\epsppr/4$.
No projected-gradient repair or additional row scan is required.
\end{corollary}
\begin{proof}
Put $\vu=\vx_\rho^*$ and $\ve=\vu-\widehat\vx$.
The preceding clipping argument gives $0\leq\ve\leq2\delta\vomega$.
Since $|\mQ|\vomega=\vomega$, we have
$|\mQ\ve|\leq2\delta\vomega$.
The RPPR residual lies in $[0,\alpha\rho\vomega]$, proving the upper
bound in \cref{eq:revisit-positive-residual}.
At a positive output coordinate, $u_i>0$ and the RPPR residual equals
$\alpha\rho\omega_i$, so the new residual is at least
$(\alpha\rho-2\delta)\omega_i\geq0$.
At a zero output coordinate, nonnegativity follows directly from the
Stieltjes signs and the nonnegative source.
The stated tolerance corresponds to $\delta=\alpha\epsppr/4$;
the freeze ratio is $144/(\alpha^2\epsppr)$.
Finally $\alpha^{-1}\mD^{1/2}(\vb-\mQ\widehat\vx)$ is a nonnegative
residual load bounded by $\epsppr\mD\one$, giving the usual ACL
residual representation of the sparse PPR output.
\end{proof}

\subsection{A threshold reporter with no mass search}
We give sufficient detail to exclude a hidden full-support scan.
Let $\sigma_k=\chi^k$ and maintain degree-density records
\begin{equation}
 X_{k,i}=x_{k,i}/(\sigma_k\omega_i),\qquad
 \mathcal R_{k,i}=[(\mQ-\mu\mI)\vx_k]_i/(\sigma_k\omega_i),\qquad
 \mathcal T_{k,i}=[(\mQ-\mu\mI)\vz_k]_i/\omega_i.
 \label{eq:revisit-lazy-records}
\end{equation}
Represent each exposed projection coordinate by the numerical key
\begin{equation}
 k_i=-\frac{\mathcal R_{k,i}}{\theta(1+\theta)}
 +\frac1{\sigma_k}\left(
       \frac{\chi z_{k,i}}{\omega_i}
       -\frac{\mathcal T_{k,i}}{1+\theta}
       +\frac{b_i}{\theta\omega_i}\right).
 \label{eq:revisit-key}
\end{equation}
Direct algebra gives
\begin{equation}
 \frac{q_{k,i}}{\omega_i}=\sigma_k k_i-\lambda_k/\theta.
 \label{eq:revisit-key-threshold}
\end{equation}
Consequently the next kinetic support is exactly the strict suffix
$k_i>\lambda_k/(\theta\sigma_k)$ in an ordered tree. Equalities return
zero and cause no row opening. No root search, weighted-sum tree, projection
multiplier, or complementarity margin is involved.

The ordinary part $-\mathcal R_{k,i}/(\theta(1+\theta))$ changes only
when a kinetic update touches that coordinate or a neighbor. The remaining
bracket is exceptional only at the seed, the current kinetic support, and
its boundary. Its entries can all be removed by their stored old keys and
reinserted after a step at the cost of those incidence sets. The exact
primal update is
\[
 \sigma_{k+1}=\chi\sigma_k,\qquad
 X_{k+1,i}=X_{k,i}+
                   \frac{\theta z_{k+1,i}}{\sigma_{k+1}\omega_i}.
\]
Only positive new kinetic coordinates change normalized primal records.
One original row scan per such coordinate updates its normalized response
and its new kinetic response on its neighbors and diagonal. All key updates
are completed before the next threshold query. The changing regularizer
changes a single scalar threshold, not all stored keys.

Use deterministic balanced trees for labels and keys. A strict suffix
enumeration takes $O(\log(M+2)+|\supp\vz_{k+1}|)$ operations with $M$
represented labels. Point updates take $O(\log(M+2))$. Thus one step has
cost
\[
 O\!\left((1+\vol(\supp\vz_k)+\vol(\supp\vz_{k+1}))\log(M+2)\right).
\]
The initial seed degree costs one reply. Every other represented label is
exposed from a charged kinetic row. A never-exposed label has zero source,
state, and neighbor response, hence raw density $-\lambda_k/\theta<0$.
The finite reporter therefore agrees exactly with the full-space update.
Inactive boundary labels incur scalar degree/key work but no row scan.

The union of historical primal supports is contained in the union of
kinetic supports. Its terminal materialization is charged once to that
history; final clipping needs no new row scan. The number of stored label
records and cached incidences is $O(1+\sum_k\vol(\supp\vz_{k+1}))$.
Every scalar iteration is affordable as well: $r_k\leq1$ gives
$K\leq\sum_k1/r_k$. These observations and
\cref{eq:revisit-explicit-work} prove the full work and storage assertions
of \cref{thm:revisit-semantic}.
